% =========================================================================
% p2_sec01_introduction.tex — Part II introduction (bridge to Part I).
% =========================================================================
\section{Introduction}
\label{sec:intro}

The companion Part~I~\cite{seni142part1} delivered the open structural layer of SENI-142: a measured, georeferenced 142-bus benchmark of the complete Dominican bulk transmission system; the HBCF diakoptic partition into the ten territorial microgrids of Law 368-22 with an exact Schur condensation at $\delta = 2.788 \times 10^{-14}\ \mathrm{pu}$; a four-phase fail-closed digital-twin lifecycle with the OSB/OVB release architecture and a 41/41-verified data-driven harness; and the definition of the operator's own $N-1$ catalog of 215 contingencies under a two-layer reserve doctrine and a four-clause static PASS battery. Part~I is the network and the method; this Part~II is the campaign that runs on it.

This paper contributes four operational results. First, a \textbf{multi-temporal AC optimal power flow} coupling power-to-gas hydrogen (electrolysis, 700~bar storage, fuel-cell re-injection), 764~MW of reversible pumped hydro, and ADM1-based biomethanation, stated with its full Karush--Kuhn--Tucker system and an exact-penalty proof that statically feasible states carry zero unserved energy. Second, a \textbf{silicon campaign} over 87,360 hourly states --- 520 complete consecutive weeks of the 2016--2025 operational archive --- resolving each national AC state in a measured 2.65~ms with 473~MiB of resident memory on 8-thread consumer hardware, with the fidelity of the condensed solve bounded by the residual and conditioning measurements of Part~I. Third, a \textbf{common-mode resilience analysis} of the island's most severe credible events: the simultaneous trip of the 752~MW Punta Catalina complex and of its sole 345~kV evacuation trunk, closed with zero unserved load under co-optimized reserve redispatch and a documented fast-storage discharge vector. Fourth, the \textbf{policy reading} of those results for hurricane-exposed insular systems: territorial microgrids as a computational restoration substrate, standards alignment, and data sovereignty.

Section~\ref{sec:opf} states the security framework recap and the multi-temporal AC-OPF with its multi-vector coupling. Section~\ref{sec:multivector} develops the flexibility constellation and its zonal economic validation. Section~\ref{sec:benchmarking} reports the silicon campaign, the accuracy measurements, and the contingency resilience battery. Section~\ref{sec:discussion} discusses the policy implications and the limitations, and Section~\ref{sec:conclusion} concludes Part~II.
